\documentclass[10pt,a4paper]{extarticle}

\usepackage[margin=1.4cm]{geometry}
\usepackage{enumitem}
\usepackage{titlesec}
\usepackage{xcolor}
\usepackage{hyperref}
\usepackage{tabularx}
\usepackage{array}

\definecolor{cairnblue}{HTML}{1a1a2e}
\definecolor{cairngrey}{HTML}{55595f}
\definecolor{cairnred}{HTML}{c8202f}
\hypersetup{colorlinks=true,linkcolor=cairnblue,urlcolor=cairnblue}

\setlength{\parindent}{0pt}
\setlength{\parskip}{0.5em}
\titleformat{\section}{\large\bfseries\color{cairnblue}}{}{0em}{}[{\color{cairnblue}\titlerule}]
\titlespacing*{\section}{0pt}{1.0em}{0.4em}
\setlist[itemize]{leftmargin=1.2em,itemsep=0.25em,topsep=0.2em,parsep=0pt}
\newcommand{\stat}[1]{{\Large\bfseries\color{cairnred}#1}}

\begin{document}

{\LARGE\bfseries\color{cairnblue} CAIRN}\par\vspace{0.2em}
{\large Rescue drones that find the people who can't call for help}\par\vspace{0.1em}
{\small\color{cairngrey} Larpverse \;\textbullet\; AgentVerse 2026 Robotics \;\textbullet\; Figures are simulation or model results unless marked as real data.}

\section*{What it is}
CAIRN is a team of heavy-lift drones that search a collapsed building for survivors \textbf{by detecting their breathing}. A rescue lead circles an area on a tablet. The drones fly there on their own, map the pile with \textbf{LiDAR}, land on it, and use a \textbf{micro-Doppler radar} to pick up the breathing motion of people buried underneath. They send back a ranked list of where to dig first, with the reason for each spot. It finds people who are unconscious or too weak to knock or shout, the people today's search misses.

\section*{How it works}
\begin{itemize}
  \item \textbf{Pick an area.} The rescue lead selects a pile on the dashboard. The drones leave the truck and fly to it.
  \item \textbf{Map and land (LiDAR).} A 128-channel LiDAR builds a 3D point-cloud map of the pile. The drones use it to find flat, stable spots (low slope and roughness) and land there.
  \item \textbf{Sense breathing (micro-Doppler radar).} A downward-facing 0.4--2\,GHz radar looks through up to about 3\,m of rubble. Phase demodulation of the radar's slow-time signal recovers each person's \textbf{chest-motion waveform}, millimetres of movement in the 0.1--0.75\,Hz breathing band. Rate, depth, regularity and pauses classify the pattern, and debris shifts, machinery vibration and rotor reflections are rejected as ``not a person''.
  \item \textbf{Rank who to reach first.} Gemini reads the breathing waveform together with what the microphones heard, and scores each spot 0--100 with a plain reason, e.g.\ ``irregular breathing with a 14\,s pause: check first''. Code checks every answer, and a rule-based fallback takes over if the AI is unavailable.
  \item \textbf{Adapt live.} When something new appears, such as a new voice or a new area, the plan changes on screen and the drones re-route. No one has to intervene.
  \item \textbf{Help on the spot.} Each drone carries a loudspeaker (``help is coming'') and an open Wi-Fi hotspot, so trapped people's phones can send ``I'm trapped''.
\end{itemize}

\section*{Why it matters}
\begin{tabularx}{\linewidth}{@{}>{\raggedright\arraybackslash}X>{\raggedright\arraybackslash}X>{\raggedright\arraybackslash}X@{}}
\stat{28\%} & \stat{0 after 62\,h} & \stat{$<$0.85\%} \\
of injured survivors in the 2003 Bam earthquake were \textbf{unconscious} when rescued, so they can't be heard (Nia et al.\ 2008). &
Weiguan Jinlong collapse, Tainan 2016: 164 rescued alive on day 1, 7 on day 2, 4 on day 3, then \textbf{nobody}. 115 died. &
of deaths are matched by live rescues from international teams, who arrive 12--56\,h late (Rom \& Kelman 2020). \\
\end{tabularx}

Today's search relies on dogs, listening devices and hand-carried radar, one spot at a time on an unstable pile. In 15 major collapses we reviewed, \textbf{no drone has ever located a living person inside the rubble}. CAIRN is built to do exactly that, in the first hour.

\section*{Results}
\begin{tabularx}{\linewidth}{@{}>{\raggedright\arraybackslash}X>{\raggedright\arraybackslash}X>{\raggedright\arraybackslash}X>{\raggedright\arraybackslash}X@{}}
\stat{3$\times$} & \stat{6 / 6} & \stat{86\%} & \stat{0} \\
as many survivors found as simpler search strategies with the same drones (17/60 vs 6/60 and 5/60; simulation, best case) &
real breath-holds detected on real radar recordings, with no false alarms (Schellenberger 2020) &
of breathing conditions identified correctly (simulation) &
false alarms: debris, machinery and rotors are rejected (simulation) \\
\end{tabularx}

\begin{itemize}
  \item \textbf{Fast:} the first radar ping is about \textbf{49 minutes after the quake}; national teams took up to 12 hours to arrive at Weiguan (model).
  \item \textbf{Worth it:} a city kit (3 drones with radar, plus van fit-out) costs about \$673k and \textbf{pays for itself if it saves one life every 18 years} (Taiwan's value of a statistical life; model). In a Weiguan-style collapse it could mean \textbf{about 6 more survivors} if the radar performs as designed (range 0.6--14; model).
  \item \textbf{Real system:} running in NVIDIA Isaac Sim, operated live from a dashboard, with an unscripted change during the recorded take.
\end{itemize}

{\small\textbf{Limits.} Radar results come from a simulated sensor; it hasn't been tested on real rubble yet, and the radar's frequency band needs a licence. CAIRN finds people faster; digging them out still takes as long. Every result is a place for rescuers to check, not a confirmed person.}

\vfill
{\footnotesize\color{cairngrey} Sources: Nia et al.\ 2008 (Prehosp Disaster Med) via Rom \& Kelman 2020 (BMJ Global Health); Chiu \& Tseng 2018 and Tainan City Fire Bureau (Weiguan); Schellenberger et al.\ 2020 (radar dataset); our Isaac Sim evaluation and cost model.}

\end{document}
